The dual space of a Hilbert space is very simple: it is isometric to the space itself. This result is known as the Riesz Representation Theorem and is one of the foundations of the theory of Hilbert spaces.

\begin{Theorem}{Riesz representation}
\label{thm:hilbert-riesz-representation}
The dual space of a Hilbert space is isometric to the Hilbert space itself. In particular, for every $x\in H$ the linear functional on $H$ defined by
\begin{equation}\label{eq:hilbert-riesz-inner-product-functional}
l_x(y):=(x,y)
\end{equation}
is bounded with norm $\|l_x\|_{H^*}=\|x\|_H$. Moreover, for every $l\in H^*$, there exists a unique $x\in H$ such that
\begin{equation}\label{eq:hilbert-riesz-functional-representation}
l(y)=(x,y)\quad\text{for every }y\in H,
\end{equation}
and, furthermore, $\|x\|_H=\|l\|_{H^*}$.
\end{Theorem}

\begin{Proof}
As the statement of the theorem indicates, for every $x\in H$ we can define a linear functional $l_x$ using \eqref{eq:hilbert-riesz-inner-product-functional}. It follows from the Cauchy-Schwarz inequality that $l_x\in H^*$ and that $\|l_x\|\leq\|x\|$. Furthermore, by setting $y=x$, we find $\|l_x\|=\|x\|$.

On the other hand, let $l$ be a nonzero linear functional on $H$. We define $M=\{y\in H\mid l(y)=0\}$. Then $M$ is a closed subspace of $H$. Furthermore, $M^\perp$ is one-dimensional since for any $u$ and $v$ in $M^\perp$ we have $l(v)u-l(u)v\in M$. (This implies $l(v)u-l(u)v=0$.)

Now, let $\widetilde x$ be a unit vector in $M^\perp$. For every $y\in H$, we have the decomposition $y=(\widetilde x,y)\widetilde x+z$, where $z\in M$, and hence $l(y)=(\widetilde x,y)l(\widetilde x)$. Thus, if we let
\begin{equation}\label{eq:riesz-representing-vector-from-kernel-complement}
x=\overline{l(\widetilde x)}\widetilde x,
\end{equation}
we see that $l=l_x:=(x,\cdot)$. The equality of the norms has already been established.
\end{Proof}

\begin{Remark}
\label{rem:riesz-mappings}
In the previous proof we constructed isomorphic mappings $A_H:H^*\to H$ and $A_H^{-1}:H\to H^*$ between a Hilbert space and the corresponding element of its dual; these are called the \emph{Riesz mappings}.
\end{Remark}

\Needspace{4\baselineskip}
\begin{Example}
\label{ex:canonical-bidual-isometry}
Let $X$ be a Banach space. Then every $x\in X$ defines a linear functional on $X^*$ by the correspondence $l_xy=y(x)$. It is clear that $\|l_x\|\leq\|x\|$, and it follows from the Hahn-Banach theorem (Theorem~\ref{thm:hahn-banach-extension}), stated in the next subsection~\ref{subsec:hahn-banach-theorem}, that there is a $y\in X^*$ with $\|y\|=1$ and $y(x)=\|x\|$. This implies that $\|l_x\|=\|x\|$, and hence we have an isometry from $X$ to a subspace of $X^{**}$. It is important to characterize those Banach spaces for which this isometry is onto.
\end{Example}

\begin{Definition}
\label{def:reflexive-banach-space}
A Banach space is called \textbf{reflexive} if the isometry $x\mapsto l_x$ from $X\to X^{**}$ defined by $l_xy=y(x)$ is surjective.
\end{Definition}

It is clear from Theorem~\ref{thm:hilbert-riesz-representation} that all Hilbert spaces are reflexive. Also, the spaces $L^p(\Omega)$, $1<p<\infty$ are reflexive. By contrast, it can be shown that $L^1(\Omega)$ and $L^\infty(\Omega)$ are not reflexive. Reflexive spaces have many useful properties that make them advantageous for analysis. In the remainder of the book, we shall mostly work in Hilbert spaces, and we shall not pursue reflexive Banach spaces further. The reader is referred to the literature, e.g., \cite{DS}.
